\begin{lemma}[Lagrangian dual problems and weak duality]
  \label{lem-weak-duality}
  Consider the feasible measure programs \eqref{eq-measure-upper} and \eqref{eq-measure-lower}, with nonempty Borel support $\Xi$ and a Borel measurable surrogate $\hat g$. Write $\bm u=(u_1,\ldots,u_K)^\top$ and $\bm l=(l_1,\ldots,l_K)^\top$. The upper and lower rank constraints have nonnegative Lagrange multipliers $\bm\lambda,\bm\gamma\in\mathbb R_+^K$, respectively. The input marginal equality has a multiplier $\varphi\in\mathcal B_b(\mathcal X)$, where $\mathcal B_b(\mathcal X)$ is the space of bounded Borel measurable real-valued functions on $\mathcal X$; $\varphi$ is unrestricted in sign.

  For fixed $\bm\lambda,\bm\gamma$, the optimal marginal multipliers for the upper and lower duals are the pointwise envelopes
  \begin{align}
    \varphi^+_{\bm\lambda,\bm\gamma}(\bm x)
    &:=\sup_{\xi\in\Xi}\left\{
      \begin{aligned}
        &\mathbf1_{\{\xi\le-\hat g(\bm x)\}}
        -\sum_{k=1}^K\lambda_k\mathbf1_{\{\xi<\xi_{(k)}^{\mathrm{te}}\}}\\
        &\quad+\sum_{k=1}^K\gamma_k\mathbf1_{\{\xi\le\xi_{(k)}^{\mathrm{te}}\}}
      \end{aligned}\right\},
    \label{eq-upper-marginal-envelope}\\
    \varphi^-_{\bm\lambda,\bm\gamma}(\bm x)
    &:=\inf_{\xi\in\Xi}\left\{
      \begin{aligned}
        &\mathbf1_{\{\xi\le-\hat g(\bm x)\}}
        +\sum_{k=1}^K\lambda_k\mathbf1_{\{\xi<\xi_{(k)}^{\mathrm{te}}\}}\\
        &\quad-\sum_{k=1}^K\gamma_k\mathbf1_{\{\xi\le\xi_{(k)}^{\mathrm{te}}\}}
      \end{aligned}\right\}.
    \label{eq-lower-marginal-envelope}
  \end{align}
  Both envelopes belong to $\mathcal B_b(\mathcal X)$. Eliminating $\varphi$ gives the Lagrangian dual problems
  \begin{align}
    \overline d_\theta
    &:=\inf_{\bm\lambda,\bm\gamma\in\mathbb R_+^K}
      \left\{\bm\lambda^\top\bm u-\bm\gamma^\top\bm l
      +\mathbb E_{\mathbb P}\left[\varphi^+_{\bm\lambda,\bm\gamma}(\bm X)\right]\right\},
    \label{eq-upper-measure-dual}\\
    \underline d_\theta
    &:=\sup_{\bm\lambda,\bm\gamma\in\mathbb R_+^K}
      \left\{-\bm\lambda^\top\bm u+\bm\gamma^\top\bm l
      +\mathbb E_{\mathbb P}\left[\varphi^-_{\bm\lambda,\bm\gamma}(\bm X)\right]\right\}.
    \label{eq-lower-measure-dual}
  \end{align}
  Each problem has $2K$ nonnegative scalar decision variables. The upper objective is convex, and the lower objective is concave. Their optimal values satisfy
  \begin{equation}
    \underline d_\theta\le\underline p_{\hat f}
    \le\overline p_{\hat f}\le\overline d_\theta.
    \label{eq-weak-duality}
  \end{equation}
\end{lemma}
\begin{proof}
  The input marginal equality is equivalent to
  \[
    \int_{\mathcal X\times\Xi}\varphi(\bm x)\,\mathrm d\Pi
    =\mathbb E_{\mathbb P}\left[\varphi(\bm X)\right],
    \qquad\forall\varphi\in\mathcal B_b(\mathcal X).
  \]
  Pairing this equality with the unrestricted multiplier $\varphi$, and assigning nonnegative multipliers to the rank inequalities, gives the upper Lagrangian \citep[Section~5.1]{B.V04}:
  \begin{equation}
    \begin{aligned}
      \mathcal L^+\left(\Pi,\varphi,\bm\lambda,\bm\gamma\right)
      ={}&\int_{\mathcal X\times\Xi}
        \mathbf1_{\{\xi\le-\hat g(\bm x)\}}\,\mathrm d\Pi\\
      &+\mathbb E_{\mathbb P}\left[\varphi(\bm X)\right]
        -\int_{\mathcal X\times\Xi}\varphi(\bm x)\,\mathrm d\Pi\\
      &+\sum_{k=1}^K\lambda_k\left(
        u_k-\int_{\mathcal X\times\Xi}
        \mathbf1_{\{\xi<\xi_{(k)}^{\mathrm{te}}\}}\,\mathrm d\Pi\right)\\
      &+\sum_{k=1}^K\gamma_k\left(
        \int_{\mathcal X\times\Xi}
        \mathbf1_{\{\xi\le\xi_{(k)}^{\mathrm{te}}\}}\,\mathrm d\Pi-l_k\right).
    \end{aligned}
    \label{eq-upper-lagrangian-expanded}
  \end{equation}
  Collecting the terms involving $\Pi$ yields
  \begin{equation}
    \begin{aligned}
      \mathcal L^+={}&\mathbb E_{\mathbb P}\left[\varphi(\bm X)\right]
        +\bm\lambda^\top\bm u-\bm\gamma^\top\bm l\\
      &+\int_{\mathcal X\times\Xi}\left[
        \begin{aligned}
          &\mathbf1_{\{\xi\le-\hat g(\bm x)\}}-\varphi(\bm x)\\
          &-\sum_{k=1}^K\lambda_k\mathbf1_{\{\xi<\xi_{(k)}^{\mathrm{te}}\}}
          +\sum_{k=1}^K\gamma_k\mathbf1_{\{\xi\le\xi_{(k)}^{\mathrm{te}}\}}
        \end{aligned}\right]\,\mathrm d\Pi.
    \end{aligned}
    \label{eq-upper-lagrangian}
  \end{equation}
  Once the marginal constraint is relaxed, $\Pi$ ranges over arbitrary finite nonnegative measures, with no prescribed total mass. If the integrand is positive at any point, placing mass $M$ there and letting $M\to\infty$ makes the supremum infinite. Otherwise its integral has supremum zero, attained by the zero measure. Nonpositivity of the integrand is exactly the pointwise requirement
  $\varphi(\bm x)\ge\varphi^+_{\bm\lambda,\bm\gamma}(\bm x)$ for every $\bm x\in\mathcal X$. Consequently the upper dual function is
  \begin{equation}
    \sup_{\Pi\in\mathcal M_+(\mathcal X\times\Xi)}\mathcal L^+
    =\begin{cases}
      \mathbb E_{\mathbb P}\left[\varphi(\bm X)\right]
        +\bm\lambda^\top\bm u-\bm\gamma^\top\bm l,
        &\text{if }\varphi\ge\varphi^+_{\bm\lambda,\bm\gamma}\text{ on }\mathcal X,\\
      +\infty,&\text{otherwise}.
    \end{cases}
    \label{eq-upper-dual-function}
  \end{equation}

  For the lower minimization problem, reverse both rank-penalty signs in \eqref{eq-upper-lagrangian-expanded}, retaining the same unrestricted marginal pairing. This gives
  \begin{equation}
    \begin{aligned}
      \mathcal L^-={}&\mathbb E_{\mathbb P}\left[\varphi(\bm X)\right]
        -\bm\lambda^\top\bm u+\bm\gamma^\top\bm l\\
      &+\int_{\mathcal X\times\Xi}\left[
        \begin{aligned}
          &\mathbf1_{\{\xi\le-\hat g(\bm x)\}}-\varphi(\bm x)\\
          &+\sum_{k=1}^K\lambda_k\mathbf1_{\{\xi<\xi_{(k)}^{\mathrm{te}}\}}
          -\sum_{k=1}^K\gamma_k\mathbf1_{\{\xi\le\xi_{(k)}^{\mathrm{te}}\}}
        \end{aligned}\right]\,\mathrm d\Pi.
    \end{aligned}
    \label{eq-lower-lagrangian}
  \end{equation}
  The infimum over nonnegative measures is finite exactly when the integrand is nonnegative everywhere, or equivalently
  $\varphi(\bm x)\le\varphi^-_{\bm\lambda,\bm\gamma}(\bm x)$ for every $\bm x\in\mathcal X$. A negative integrand at a point allows an arbitrarily large point mass to drive the infimum to $-\infty$. Hence
  \begin{equation}
    \inf_{\Pi\in\mathcal M_+(\mathcal X\times\Xi)}\mathcal L^-
    =\begin{cases}
      \mathbb E_{\mathbb P}\left[\varphi(\bm X)\right]
        -\bm\lambda^\top\bm u+\bm\gamma^\top\bm l,
        &\text{if }\varphi\le\varphi^-_{\bm\lambda,\bm\gamma}\text{ on }\mathcal X,\\
      -\infty,&\text{otherwise}.
    \end{cases}
    \label{eq-lower-dual-function}
  \end{equation}

  To justify the pointwise choices, fix $\bm\lambda,\bm\gamma$ and replace $-\hat g(\bm x)$ by a real threshold $t$ in \eqref{eq-upper-marginal-envelope} and \eqref{eq-lower-marginal-envelope}. Every expression inside the supremum or infimum is nondecreasing in $t$. Both envelopes are therefore nondecreasing functions of $t$, and hence Borel measurable. They are bounded because the failure indicator lies in $[0,1]$ and the multiplier sums are finite. Composition with the Borel measurable map $-\hat g$ proves that $\varphi^+_{\bm\lambda,\bm\gamma}$ and $\varphi^-_{\bm\lambda,\bm\gamma}$ belong to $\mathcal B_b(\mathcal X)$.

  For fixed rank multipliers, the upper dual minimizes $\mathbb E_{\mathbb P}[\varphi(\bm X)]$ over functions above the upper envelope. Its minimum is attained by $\varphi=\varphi^+_{\bm\lambda,\bm\gamma}$. The lower dual maximizes the same expectation over functions below the lower envelope, with maximum attained by $\varphi=\varphi^-_{\bm\lambda,\bm\gamma}$. Substitution gives \eqref{eq-upper-measure-dual} and \eqref{eq-lower-measure-dual}. These choices require no maximizer or minimizer in the residual variable. The upper envelope is a supremum of affine functions of the rank multipliers and is therefore convex; the lower envelope is an infimum of affine functions and is concave. Expectations preserve these properties.

  Finally, for every primal feasible $\Pi$, the penalty signs and marginal equality imply
  $\mathcal L^-\le p_{\hat f}(\Pi)\le\mathcal L^+$.
  Taking the respective measure extrema and optimizing the dual variables proves \eqref{eq-weak-duality}. All steps remain valid at $\theta=0$.
\end{proof}

The functional multiplier has now been eliminated, but the dual objectives still involve an expectation of a residual supremum or infimum. To evaluate these quantities through the residual ranks, collect the rank penalties in the step function
\begin{equation}
  \eta_{\bm\lambda,\bm\gamma}(\xi)
  :=\sum_{k=1}^K\lambda_k\mathbf1_{\{\xi<\xi_{(k)}^{\mathrm{te}}\}}
    -\sum_{k=1}^K\gamma_k\mathbf1_{\{\xi\le\xi_{(k)}^{\mathrm{te}}\}}.
  \label{eq-dual-step-function}
\end{equation}
The two envelopes maximize $\mathbf1_{\{\xi\le-\hat g(\bm x)\}}-\eta_{\bm\lambda,\bm\gamma}(\xi)$ and minimize $\mathbf1_{\{\xi\le-\hat g(\bm x)\}}+\eta_{\bm\lambda,\bm\gamma}(\xi)$ over $\Xi$, respectively. The following finite aggregation makes the ordered rank structure explicit.

